\subsection{Sprint 1}

Na Sprint 1, atuei na organização do fluxo de trabalho do time e no
acompanhamento da integração das entregas. Inicialmente, trabalhei na integração
entre o ClickUp e o GitHub, buscando estabelecer padrões para a criação de
\textit{branches} e automatizar a atualização das tarefas. Para essa atividade,
contei com o apoio do colega AGES III Leonardo Dondoni.

Também prestei apoio pontual aos integrantes do time durante o desenvolvimento.
Auxiliei o Arthur Kasprzak, AGES I, com a \textit{branch} em que ele trabalhava no
\textit{backend}. Além disso, investiguei um erro relacionado às ferramentas do
projeto, pois o \textit{mirror} dos repositórios não estava funcionando. Naquele
momento, havia a suspeita de que o problema estivesse no servidor da AGES, sem que
fosse possível confirmá-la a partir da investigação realizada.

Outra dificuldade percebida no início da sprint foi a falta de um alinhamento
técnico coletivo sobre a arquitetura, a estrutura dos repositórios e a preparação
dos ambientes de desenvolvimento. Embora um dos AGES III tenha criado
os \textit{boilerplates}, as decisões que orientaram essa estrutura não foram
apresentadas ao time de forma suficiente, especialmente no que se referia ao
\textit{frontend}. Isso dificultou que os demais integrantes encontrassem
referências técnicas para começar e conduzir suas tarefas.

Na segunda metade da sprint, concentrei-me no controle das entregas que ainda
precisavam ser integradas. Ao revisar os \textit{pull requests} do
\textit{frontend}, identifiquei diversos PRs abertos havia mais de uma semana,
formando uma fila que dificultava a validação do trabalho. Nos dois últimos dias
registrados da sprint, revisei e integrei com todos os AGES IV os PRs pendentes, testei a aplicação e
busquei finalizar as histórias de usuário que permaneciam abertas. Esse esforço
teve como objetivo preparar a entrega da Sprint 1 para a apresentação às
\textit{stakeholders}.

\medskip

\fbox{%
    \parbox{0.95\linewidth}{%
        \setlength{\parindent}{15pt}
        \itshape Na minha avaliação, boa parte das dificuldades que se tornaram
        evidentes nesta sprint teve origem na desorganização inicial de nós, AGES
        IV. O atraso dos protótipos, concluídos somente no dia anterior à entrega
        da Sprint 0, reduziu o tempo e a clareza que o time teria para iniciar o
        desenvolvimento. Eu acredito que isso impactou principalmente os colegas
        de AGES I e II, que ainda precisavam de direcionamento para entender o
        produto e começar a trabalhar.

        Também senti falta de um momento em que as decisões técnicas fossem
        apresentadas ao grupo e os ambientes de desenvolvimento fossem preparados
        coletivamente. Um dos AGES III gerou os \textit{boilerplates}, mas não
        houve uma explicação suficiente sobre os motivos da organização adotada.
        Como esse colega tinha maior foco no \textit{backend}, a estrutura do
        \textit{frontend} ficou especialmente sem referência. Ao meu ver, isso
        contribuiu para que o time encontrasse dificuldades para avançar e para o
        acúmulo de PRs que encontrei abertos ao final da sprint.

        Não considero que a única causa dos problemas tenha sido essa falta de
        alinhamento, mas ela me pareceu uma das mais relevantes. Esperava que os
        AGES III atuassem como referências técnicas, ajudando os iniciantes a
        compreender a arquitetura e revisando as entregas desde cedo. Em vez disso,
        percebi que assumiram tarefas de desenvolvimento logo no início, o que, na
        minha visão, reduziu oportunidades de aprendizado para os AGES I e II sem
        estabelecer as referências de que o time precisava. Essa situação foi
        discutida entre nós, AGES IV, e também levantada na retrospectiva, reforçando
        a importância de definir responsabilidades e direcionar o time antes de
        acelerar a implementação.
    }%
}
